\documentclass[11pt,a4paper]{article}
\usepackage[utf8]{inputenc}
\usepackage[T1]{fontenc}
\usepackage[margin=0.72in]{geometry}
\usepackage{enumitem}
\usepackage{titlesec}
\usepackage{xcolor}
\usepackage{hyperref}
\usepackage{microtype}

\definecolor{primarycolor}{RGB}{33,37,41}
\definecolor{secondarycolor}{RGB}{0,91,150}
\definecolor{accentcolor}{RGB}{95,105,115}
\hypersetup{colorlinks=true,urlcolor=secondarycolor,linkcolor=secondarycolor}
\titleformat{\section}{\large\bfseries\color{primarycolor}}{}{0em}{}
\titlespacing*{\section}{0pt}{1.0em}{0.35em}
\setlist[itemize]{leftmargin=1.4em,itemsep=0.15em,topsep=0.2em,parsep=0em}
\pagestyle{empty}
\setlength{\parindent}{0pt}
\setlength{\parskip}{0.25em}

\begin{document}

\begin{center}
{\LARGE\bfseries MARTA KOWAL-BRANDT}\\
QA Automation Engineer | DevOps / Platform Engineering Transition\\
Wrocław, Poland | marta.kowalbrandt@example.com | +48 555 303 303
\end{center}

\section*{PROFESSIONAL SUMMARY}
QA automation engineer with 7 years of experience, including two years focused largely on CI/CD pipelines, test environments and build tooling. Owns GitLab CI pipelines and namespace-level ephemeral Kubernetes test environments for three product teams, improved average pipeline time from 38 to 14 minutes and reduced flaky tests from 7\% to below 1\%, and is seeking a DevOps or platform engineering role.

\section*{LANGUAGES}
Polish (native), English (C1), German (A2)

\section*{TECHNICAL TOOLKIT}
CI/CD and containers: GitLab CI, Jenkins, Docker, Docker Compose, Helm\\
Infrastructure and cloud: Kubernetes (namespace-level access at work; k3s personal homelab), Terraform (personal project), Ansible (basic), AWS (learning: EC2 and S3 basics)\\
Monitoring and data: Prometheus, Grafana, Loki, PostgreSQL, SQL\\
Programming and scripting: Python, Bash, Java (test code)\\
Testing: pytest, Playwright, Selenium, Postman\\
Systems and tools: Linux, Git, Argo CD, Traefik, cert-manager

\section*{PROFESSIONAL EXPERIENCE}
\textbf{QA Automation Engineer (Test Infrastructure) | Odrava Health Software, Wrocław} \hfill \textit{2021–Present}
\begin{itemize}
\item Own GitLab CI pipelines for three product teams of approximately 25 developers, covering build, unit, integration and end-to-end test stages, and artifact publishing. Improved average pipeline time from 38 to 14 minutes by parallelising test jobs, caching dependencies and separating end-to-end tests into a dedicated stage; monthly CI runner costs fell by about 30\%.
\item Built merge-request test environments using Docker Compose and later Helm charts deployed to a shared Kubernetes cluster; environments include the application, PostgreSQL and seeded data.
\item Maintain a pytest and Playwright end-to-end suite of approximately 900 tests; reduced the flaky test rate from 7\% to below 1\%.
\item Wrote Bash and Python tooling to seed PostgreSQL test data, and created Grafana dashboards for flaky-test metrics using CI results and PostgreSQL. Wrote an internal guide to the CI system used in onboarding and mentored two QA interns in 2023.
\end{itemize}

\textbf{QA Engineer | Brightmill Software House, Wrocław} \hfill \textit{2019–2021}
\begin{itemize}
\item Tested web and mobile applications through manual and automated testing for client projects.
\item Introduced a Selenium and Java automated regression suite for the largest client project and configured Jenkins jobs for nightly runs.
\end{itemize}

\section*{SELECTED INFRASTRUCTURE PROJECTS}
\textbf{Homelab k3s Cluster with GitOps | Personal} \hfill \textit{2024–2025}
\begin{itemize}
\item Set up a three-node Raspberry Pi 5 k3s cluster, with applications deployed from a Git repository using Argo CD.
\item Added Prometheus, Grafana and Loki for monitoring, Traefik ingress, and cert-manager for TLS.
\item Provisioned infrastructure for a small Hetzner VPS using Terraform; used Ansible at a basic level.
\end{itemize}

\textbf{Merge-Request Ephemeral Environments | Odrava Health} \hfill \textit{2023–2024}
\begin{itemize}
\item Implemented short-lived environments for merge requests, deployed through GitLab CI and Helm to the shared Kubernetes cluster.
\item Configured automatic teardown after merge or 48 hours of inactivity; QA and product teams use the environments to review features, and end-to-end tests run against them.
\end{itemize}

\section*{EDUCATION \& TRAINING}
BEng Computer Engineering | Silesian Lowlands University of Technology | 2018\\
ISTQB Foundation Level | 2019\\
Certified Kubernetes Administrator (CKA): studying; exam booked for January 2027

\end{document}